\section{Exact reduction and sharp spectral asymptotics}
\label{sec:exact-spectrum}

The preceding argument controls every polymer observable without assuming
that the full gap is attained by a count observable. We now prove that it
is attained there. Granovsky and Erlihson
\cite[Section~2.3]{GranovskyErlihson2009} already announce this equality
for their solvable class. The purpose of the argument below is an
independent proof which explicitly controls all the orthogonal sectors.
It also supplies the comparison needed for the spatial extension.

\begin{lemma}[Adjacent ratios]\label{lem:ground-ratio}
Let $h>0$ be a ground state of the count-derivative operator
$\mathcal A$ in~\eqref{proof:eq:derivativegenerator}:
$-\mathcal Ah=\gamma_Qh$. Then
\[
 \frac{h(i+1)}{h(i)}>\frac12,\qquad 1\le i\le n-2.
\]
For $n=2$ the assertion has no ratios to check.
\end{lemma}

\begin{proof}
Write $\lambda_i=\beta(n-i-1)$, $\mu_i=i(i-1)/n$ and
$V_i=\beta+2i/n$. Suppose that $r=h(i+1)/h(i)\le1/2$ is the
minimum of the forward ratios. The eigenvalue equations at $i$ and
$i+1$, respectively, imply
\begin{align*}
 \gamma_Q&\ge\lambda_i(1-r)+\mu_i(1-r^{-1})+V_i,\\
 \gamma_Q&\le\lambda_{i+1}(1-r)+\mu_{i+1}(1-r^{-1})+V_{i+1}.
\end{align*}
Indeed, the preceding ratio is at least $r$ in the first inequality,
and the following ratio is at least $r$ in the second. A missing
preceding or following ratio is multiplied by a zero boundary rate.
Subtracting gives
\[
 0\ge\beta(1-r)+\frac{2i}{n}(r^{-1}-1)-\frac2n>0.
\]
At $r=1/2$ the last expression equals $\beta/2+2(i-1)/n>0$;
it increases as $r$ decreases. This contradiction
proves the assertion.
\end{proof}

\begin{proposition}[Comparison of the higher sectors]
\label{prop:higher-exact}
For $2\le\ell\le\lfloor n/2\rfloor$,
\[
 \inf\operatorname{Spec}(-\mathcal L_\ell)
 \ge\gamma_Q+\frac{(\ell-2)(\ell+1)}n.
\]
In fact every sector with $\ell\ge2$ has its bottom eigenvalue
strictly above $\gamma_Q$.
\end{proposition}

\begin{proof}
Set $r=k-\ell+1$ and use $g(k)=h(r)$ as a positive test vector
for~\eqref{proof:eq:sector}. Its domain is
$1\le r\le n-2\ell+1$. Comparing the birth, death and killing
coefficients with those of $\mathcal A$ gives
\begin{align}
 -\frac{\mathcal L_\ell g(k)}{g(k)}
 ={}&\gamma_Q
 +(\ell-1)\beta\left(2\frac{h(r+1)}{h(r)}-1\right)\notag\\
 &+\frac{\ell-2}{n}
 \left[r+\ell-(r-1)\left(\frac{h(r-1)}{h(r)}-1\right)\right].
 \label{eq:higher-exact-certificate}
\end{align}
The forward-ratio term is strictly positive by
Lemma~\ref{lem:ground-ratio}. The backward ratio is at most two,
so the expression in square brackets is at least $\ell+1$.
The positive-vector comparison proves the claim. At $r=1$ the
backward term is zero. At the upper endpoint,
$r+1=n-2\ell+2\le n-1$, so the forward value is in the domain
of $h$. The same computation applies to a one-state sector.
\end{proof}

\begin{theorem}[Exact count reduction]\label{thm:exact-gap}
For every $n\ge2$ and $\beta>0$,
\[
 \boxed{\gamma_M(n,\beta)=\gamma_Q(n,\beta).}
\]
More generally, the best Poincar\'e constant on rotation-invariant
cut-set observables is $\gamma_Q$.
\end{theorem}

\begin{proof}
Rotation invariance eliminates degree one, as proved in
Proposition~\ref{proof:prop:base}. The nonconstant degree-zero
sector has lowest rate $\gamma_Q$, and
Proposition~\ref{prop:higher-exact} puts every surviving higher
sector at or above that rate. Orthogonality therefore gives a
Poincar\'e inequality with constant $\gamma_Q$ on all
rotation-invariant functions. The stationary comparison transfers
it to polymers. The opposite inequality follows by testing with
count observables. For $n=2,3$ there is no higher sector to check.
\end{proof}

\subsection{An explicit error bound for the leading constant}

The same family of positive vectors used for the earlier upper
bound gives a sharper result if its parameter is chosen by an
algebraic cancellation.

\begin{theorem}[Sharp estimate with a uniform error]
\label{thm:sharp-gap}
Put
\[
 d=\sqrt{\beta^2+4\beta},\qquad
 a=\frac{\beta n+1+
       \sqrt{(\beta n+1)^2+4\beta n(n+1)}}2,
 \qquad c=\frac{2a}{n}-\beta.
\]
Then $a>1$ and
\begin{equation}\label{eq:sharp-bracket}
 c-\frac{6(a-1)}{n(a+3)}
 \le\gamma_M
 \le c-\frac{6a(a-1)}{n(n+a-2)(n+a+1)}.
\end{equation}
In particular,
\begin{equation}\label{eq:sharp-additive}
 \boxed{\quad d-\frac5n\le\gamma_M(n,\beta)
                     \le d+\frac4n.\quad}
\end{equation}
These estimates can be combined with
$\gamma_M\ge\beta+2/n$ whenever the lower expression is small.
\end{theorem}

\begin{proof}
Choose
$h(k)=[(k+a)(k+a+1)]^{-1}$ for the count derivative.
Its forward and backward ratios are respectively
$1-2/(k+a+2)$ and $1+2/(k+a-1)$. Hence
\begin{align*}
 -\frac{\mathcal Ah(k)}{h(k)}
 &=\beta+\frac{2\beta(n-k-1)}{k+a+2}
                       +\frac{2ak}{n(k+a-1)}\\
 &=\frac{2a}{n}-\beta+
   \frac{2\beta(n+a+1)}{k+a+2}
             -\frac{2a(a-1)}{n(k+a-1)}.
\end{align*}
The definition of $a$ is precisely
$a(a-1)=\beta n(n+a+1)$. Thus
\begin{equation}\label{eq:sharp-certificate}
 -\frac{\mathcal Ah(k)}{h(k)}
 =c-\frac{6a(a-1)}
          {n(k+a-1)(k+a+2)}.
\end{equation}
This quantity increases with $k$; its extreme values at $1$ and
$n-1$ give~\eqref{eq:sharp-bracket} by positive-vector comparison
and Theorem~\ref{thm:exact-gap}. Both endpoint terms in the original
generator are treated with their zero coefficients, so the calculation
also holds for $n=2$.

The subtracted term in~\eqref{eq:sharp-certificate} is between
zero and $6/n$. Moreover,
\[
 c=\frac1n+\sqrt{d^2+\frac{6\beta}{n}+\frac1{n^2}},
 \qquad d+\frac1n\le c\le d+\frac4n.
\]
For the upper inequality, square $d+3/n$ and use $d\ge\beta$.
The lower inequality is immediate. These facts imply
\eqref{eq:sharp-additive}.
\end{proof}

\begin{corollary}[Asymptotic equivalents]\label{cor:sharp-equivalents}
Let $n\to\infty$ and allow $\beta=\beta_n>0$ to vary.
If $n\sqrt{\beta_n}\to\infty$, then
\[
 \gamma_M(n,\beta_n)\sim\sqrt{\beta_n^2+4\beta_n}.
\]
In particular,
\begin{align*}
 \beta_n\sim b/n,\ b>0
 &\quad\Longrightarrow\quad
 \sqrt n\,\gamma_M(n,\beta_n)\longrightarrow2\sqrt b,\\
 \beta_n\longrightarrow\beta_0>0
 &\quad\Longrightarrow\quad
 \gamma_M(n,\beta_n)\longrightarrow\sqrt{\beta_0^2+4\beta_0}.
\end{align*}
If $n^2\beta_n\to0$, then
\[
 n\gamma_M(n,\beta_n)\longrightarrow2.
\]
\end{corollary}

\begin{proof}
In the first regime, $nd\ge2n\sqrt{\beta_n}\to\infty$,
so the additive error in~\eqref{eq:sharp-additive} is negligible.
The two displayed specializations follow. In the final regime,
$a\to1$ in Theorem~\ref{thm:sharp-gap}; its upper bound gives
$n\gamma_M\le2a-\beta_n n\to2$, while the lower bound
$\gamma_M\ge\beta_n+2/n$ gives the reverse limit.
\end{proof}
